\appendix
\section{Additional Derived-Pointer Examples}\label{app:derived_pointers}

This appendix presents additional examples of object-owned memory and subordinate pointers in CRuby, supplementing the string example in Section~\ref{sec:Background}.

\subsection*{Arrays}
Arrays similarly have an object header plus an element store.  CRuby may embed a small
number of elements directly in the object, or it may use a separate buffer referenced
by a pointer field (conceptually \texttt{ptr} in \texttt{struct RArray}).\footnote{\url{https://docs.ruby-lang.org/capi/en/master/dd/d8b/struct_r_array.html}
(describes embedded elements vs.\ heap-backed arrays and the \texttt{ptr} field).}

When code needs a raw \texttt{VALUE *} pointing to the contiguous element buffer,
CRuby provides APIs and macros such as \texttt{RARRAY\_CONST\_PTR} and the scoped
\texttt{RARRAY\_PTR\_USE} wrapper.  \texttt{RARRAY\_PTR\_USE} makes the ``raw pointer
region'' explicit and delegates setup/teardown to internal helpers (Listing~\ref{lst:rarray_ptr_use}).
This scoping is important in practice because arrays can share storage and can change
representation as they grow or as operations demand copies; treating a raw
\texttt{VALUE *} as a long-lived stable pointer is therefore brittle.

\begin{figure}[tb]
\begin{lstlisting}[caption={Scoped raw-pointer access for arrays via \texttt{RARRAY\_PTR\_USE} (from \texttt{rarray.h} documentation).},label={lst:rarray_ptr_use}]
const VALUE ary = rb_eval_string("[...]");
const int len = RARRAY_LENINT(ary);

RARRAY_PTR_USE(ary, ptr, {
  rb_funcallv(rb_stdout, rb_intern("write"), len, ptr);
});
\end{lstlisting}
\end{figure}

\subsection*{Typed Data}
Beyond core containers, a common ``object-owned memory'' pattern is \texttt{T\_DATA}
(typed data), where a Ruby object owns a pointer to an extension-defined C struct.
The ownership contract is described by an \texttt{rb\_data\_type\_t} that supplies
callbacks (notably \texttt{dfree}) that the GC invokes when the wrapper object is
collected.  The pointer is extracted with \texttt{TypedData\_Get\_Struct}.\footnote{\url{https://docs.ruby-lang.org/en/3.1/extension_rdoc.html} (sections
``Encapsulate C Data into a Ruby Object'' and \texttt{TypedData\_Get\_Struct}).}
From the perspective of this paper, the extracted \texttt{struct *} is a derived
pointer exactly like \texttt{RSTRING\_PTR}: it is valid only while the owning
\texttt{VALUE} remains live.

\section{Source Code of \texttt{rb\_marshal\_dump\_limited}}\label{app:rb_marshal_dump_limited}
Quoted from CRuby v3.4.5.\footnote{\url{https://github.com/ruby/ruby/blob/20cda200d3ce092571d0b5d342dadca69636cb0f/marshal.c#L1211-L1248}}

\begin{figure}[tb]
\begin{lstlisting}[caption={Source code of \texttt{rb\_marshal\_dump\_limited}},label={lst:rb_marshal_dump_limited_full}]
VALUE
rb_marshal_dump_limited(VALUE obj, VALUE port, int limit)
{
    struct dump_arg *arg;
    VALUE wrapper; /* used to avoid memory leak in case of exception */

    wrapper = TypedData_Make_Struct(0, struct dump_arg, &dump_arg_data, arg);
    arg->dest = 0;
    arg->symbols = st_init_numtable();
    arg->data    = rb_init_identtable();
    arg->num_entries = 0;
    arg->compat_tbl = 0;
    arg->encodings = 0;
    arg->userdefs = 0;
    arg->str = rb_str_buf_new(0);
    if (!NIL_P(port)) {
        if (!rb_respond_to(port, s_write)) {
            io_needed();
        }
        arg->dest = port;
        dump_check_funcall(arg, port, s_binmode, 0, 0);
    }
    else {
        port = arg->str;
    }

    w_byte(MARSHAL_MAJOR, arg);
    w_byte(MARSHAL_MINOR, arg);

    w_object(obj, arg, limit);
    if (arg->dest) {
        rb_io_write(arg->dest, arg->str);
        rb_str_resize(arg->str, 0);
    }
    clear_dump_arg(arg);
    RB_GC_GUARD(wrapper);

    return port;
}
\end{lstlisting}
\end{figure}

\section{Source Code of \texttt{ossl\_cipher\_update}}\label{app:ossl_cipher_update}
Quoted from CRuby v3.4.5.\footnote{\url{https://github.com/ruby/ruby/blob/20cda200d3ce092571d0b5d342dadca69636cb0f/ext/openssl/ossl_cipher.c#L372}}

%\begin{figure}[tb]
\begin{lstlisting}[caption={Source code of \texttt{ossl\_cipher\_update}},label={lst:ossl_cipher_update_full}]
	    /*
 *  call-seq:
 *     cipher.update(data [, buffer]) -> string or buffer
 *
 *  Encrypts data in a streaming fashion. Hand consecutive blocks of data
 *  to the #update method in order to encrypt it. Returns the encrypted
 *  data chunk. When done, the output of Cipher#final should be additionally
 *  added to the result.
 *
 *  If _buffer_ is given, the encryption/decryption result will be written to
 *  it. _buffer_ will be resized automatically.
 */
static VALUE
ossl_cipher_update(int argc, VALUE *argv, VALUE self)
{
    EVP_CIPHER_CTX *ctx;
    unsigned char *in;
    long in_len, out_len;
    VALUE data, str;

    rb_scan_args(argc, argv, "11", &data, &str);

    if (!RTEST(rb_attr_get(self, id_key_set)))
	ossl_raise(eCipherError, "key not set");

    StringValue(data);
    in = (unsigned char *)RSTRING_PTR(data);
    in_len = RSTRING_LEN(data);
    GetCipher(self, ctx);

    /*
     * As of OpenSSL 3.2, there is no reliable way to determine the required
     * output buffer size for arbitrary cipher modes.
     * https://github.com/openssl/openssl/issues/22628
     *
     * in_len+block_size is usually sufficient, but AES key wrap with padding
     * ciphers require in_len+15 even though they have a block size of 8 bytes.
     *
     * Using EVP_MAX_BLOCK_LENGTH (32) as a safe upper bound for ciphers
     * currently implemented in OpenSSL, but this can change in the future.
     */
    if (in_len > LONG_MAX - EVP_MAX_BLOCK_LENGTH) {
	ossl_raise(rb_eRangeError,
		   "data too big to make output buffer: %ld bytes", in_len);
    }
    out_len = in_len + EVP_MAX_BLOCK_LENGTH;

    if (NIL_P(str)) {
        str = rb_str_new(0, out_len);
    } else {
        StringValue(str);
        if ((long)rb_str_capacity(str) >= out_len)
            rb_str_modify(str);
        else
            rb_str_modify_expand(str, out_len - RSTRING_LEN(str));
    }

    if (!ossl_cipher_update_long(ctx, (unsigned char *)RSTRING_PTR(str), &out_len, in, in_len))
	ossl_raise(eCipherError, NULL);
    assert(out_len <= RSTRING_LEN(str));
    rb_str_set_len(str, out_len);

    return str;
}
\end{lstlisting}
%\end{figure}

\section{Source code of \texttt{ruby\_brace\_expand}}\label{app:ruby_brace_expand}
Quoted from CRuby v3.4.5.\footnote{\url{https://github.com/ruby/ruby/blob/20cda200d3ce092571d0b5d342dadca69636cb0f/dir.c#L3193}}

\begin{figure}[tb]
\begin{lstlisting}[caption={Source code of \texttt{ruby\_brace\_expand}},label={lst:ruby_brace_expand_full}]
	static int
	ruby_brace_expand(const char *str, int flags, ruby_glob_func *func, VALUE arg,
	                  rb_encoding *enc, VALUE var)
{
    const int escape = !(flags & FNM_NOESCAPE);
    const char *p = str;
    const char *pend = p + strlen(p);
    const char *s = p;
    const char *lbrace = 0, *rbrace = 0;
    int nest = 0, status = 0;

    while (*p) {
        if (*p == '{' && nest++ == 0) {
            lbrace = p;
        }
        if (*p == '}' && lbrace && --nest == 0) {
            rbrace = p;
            break;
        }
        if (*p == '\\' && escape) {
            if (!*++p) break;
        }
        Inc(p, pend, enc);
    }

    if (lbrace && rbrace) {
        size_t len = strlen(s) + 1;
        char *buf = GLOB_ALLOC_N(char, len);
        long shift;

        if (!buf) return -1;
        memcpy(buf, s, lbrace-s);
        shift = (lbrace-s);
        p = lbrace;
        while (p < rbrace) {
            const char *t = ++p;
            nest = 0;
            while (p < rbrace && !(*p == ',' && nest == 0)) {
                if (*p == '{') nest++;
                if (*p == '}') nest--;
                if (*p == '\\' && escape) {
                    if (++p == rbrace) break;
                }
                Inc(p, pend, enc);
            }
            memcpy(buf+shift, t, p-t);
            strlcpy(buf+shift+(p-t), rbrace+1, len-(shift+(p-t)));
            status = ruby_brace_expand(buf, flags, func, arg, enc, var);
            if (status) break;
        }
        GLOB_FREE(buf);
    }
    else if (!lbrace && !rbrace) {
        status = glob_call_func(func, s, arg, enc);
    }

    RB_GC_GUARD(var);
    return status;
}
\end{lstlisting}
\end{figure}
